% !TeX spellcheck = es_ES
\documentclass[12pt]{article}
%\usepackage[utf8]{inputenc}
\usepackage{moodle}
\begin{document}
\begin{quiz}{Matemáticas/Funciones elementales}
% \begin{multi}[numbering=none, 
% feedback={
% \[x^2+1\]
% y hacemos Taylor
% }]{A first derivative}
%   What is the first derivative of $x^3$?
%   \item  $\frac{1}{4} x^4+C$
%   \item* $3x^2$
%   \item  $51$
% \end{multi}
 
 \begin{multi}[shuffle=true,numbering=none,
 feedback={Si llamamos $f(x)=5x+3$ tenemos que $f(x)-3=5x$ y dividiendo por $5$ se tiene que $x=\displaystyle \frac{f(x)-3}{5}.$
 }]{Pregunta funciones inversas}
 Señala cuales de los siguientes pares de funciones son inversas.
	\item*  $f(x)=5x+3$, $g(x)=\frac{x-3}{5}$
	\item $f(x)=\frac{4}{5}x+4$, $g(x)= \frac{5}{4}x+3$ 
	\item $f(x)=x^2$ si  $x<0$, $g(x)= \sqrt{x}$ si $x>0$
 \end{multi}
 
 \begin{multi}[shuffle=true,numbering=none]{Pregunta función inversa}
	 Halla la inversa de la función $f:\mathbb{R}^+_0\to \mathbb{R}$ definida por $f(x)=x^2-5$.
	\item $g(x)=\sqrt{x}+5,\ x\ge -5$ 
	\item*  $g(x)=\sqrt{x+5},\ x\ge -5$ 
 \end{multi}
 
 \begin{numerical}[shuffle=true,numbering=none,feedback={
 \[ \log_2(4)+\log_3\frac{1}{9}= \log_2(2^2)+\log_3 3^{-2}=
 2\overbrace{log_2(2)}^{1}-2\overbrace{log_3(3)}^{1}=0.\]
 }]{Pregunta expresión evaluación funciones}Calcula la siguiente expresión: $\displaystyle \log_2(4)+\log_3\frac{1}{9}$.
 
 \item 0 
 \end{numerical}
 
 \begin{numerical}[shuffle=true,numbering=none, feedback={
 \[\begin{aligned} 
 5\log_3(9)-2\log_2(16)  	& = 5\log_3(3^2)-2\log_2 (2^4) \\
 				& =10\overbrace{\log_3(3)}^{1}-8\overbrace{\log_2(2)}^{1}=0.
 \end{aligned}\]}
 ]{Pregunta sobre evaluación de funciones} 
 Simplifica la expresión  $5\log_3(9)-2\log_2(16)$.
 \item 2
 \end{numerical}
 
 \begin{numerical}[shuffle=true,numbering=none, feedback={Tenemos $1=\log_3(x)+\log_3(2x+1)=\log_3(x(2x+1))$. Si elevamos $3$ a cada uno de las expresiones primera y última nos queda
 $$3=3^1=3^{\log_3(x(2x+1))}=x(2x+1)=2x^2+x\Longrightarrow 2x^2+x-3=0$$
 Si resolvemos la ecuación de segundo grado $ 2x^2+x-3=0$ obtenemos las soluciones $x=1$ y $x=-3/2$ pero esta última solución no nos sirve ya que logaritmo  sólo se aplica a números positivos. Por tanto la solución es $x=1$.
 }]{Ecuación funciones elementales}
 Resuelve la ecuación $\log_3(x)+\log_3(2x+1)=1$ y escribe el valor de $x$.
 \item 1    
 \end{numerical}
 
 \begin{numerical}[shuffle=true,numbering=none, feedback={Se tiene que $ 2^{3\log_2(x)}=2^{\log_2(x^3)}=x^3$
 y, por otro lado, $4\log_3(9)=4\log_3(3^2)=8$  y la ecuación  a resolver queda $x^3=8$ que tiene como única solución $x=2$.
 }]
  {Ecuación funciones elementales}
   Resuelve la ecuación $ 2^{3\log_2(x)}=4\log_3(9)$ y escribe el valor de $x$.
   \item 2
 \end{numerical}

	\begin{multi}[shuffle=true,numbering=none, feedback={La fórmula del coseno del ángulo doble nos da
	$$\cos(2x)=\cos^2(x)-\mathrm{sen}^2(x)=\cos^2(x)-(1-\cos^2(x))=$$
	$$=2\cos^2(x)-1.$$
	Despejando queda $\cos^2(x)=\frac{1+\cos(2x)}{2}$.
	}]{Simplificación funciones trigonométricas}Decide a qué respuesta es igual $\cos^2(x)$.
	\item $\frac{1-\cos(2x)}{2}$  
	\item* $\frac{1+\cos(2x)}{2}$ 
	\item $1-\frac{\cos(2x)}{2}$
	\end{multi}
	
	\begin{multi}[shuffle=true,numbering=none, feedback={La fórmula del coseno del ángulo doble nos da
	$$\cos(2x)=\cos^2(x)-\mathrm{sen}^2(x)=(1-\mathrm{sen}^2(x))-\mathrm{sen}^2(x)=$$
	$$=1-2\mathrm{sen}^2(x).$$
	Despejando queda $\mathrm{sen}^2(x)=\frac{1-\cos(2x)}{2}$.
	}]{Simplificación funciones trigonométricas} Decide a qué respuesta es igual $\mathrm{sen}^2(x)$.
	\item* $\frac{1-\cos(2x)}{2}$  
	\item $\frac{1+\cos(2x)}{2}$ 
	\item $1-\frac{\cos(2x)}{2}$
	\end{multi}	
	
	\begin{multi}[shuffle=true,numbering=none, feedback={Arcoseno es la inversa de la función seno, cuando consideramos seno definida en el intervalo $[0, \pi]$. Como $0\le 2\le \pi$ entonces la igualdad es cierta.
	}]{Evaluación función trigonométrica}¿Es cierto que $\mathrm{arcsen}(\mathrm{sen}(2))=2$?
	\item* Sí 
	\item No
	\end{multi}
	
	\begin{multi}[shuffle=true,numbering=none, feedback={Arcocoseno es la inversa de la función coseno, cuando consideramos coseno definida en el intervalo $[-\pi /2, \pi /2]$. Como $ \pi /2 <2$ entonces la igualdad no es cierta.
	}]{Evaluación función trigonométrica}¿Es cierto que $\arccos(\cos(2))=2$?
	\item Sí 
	\item* No
	\end{multi}
\end{quiz}
\end{document}
